Para este sistema, se implementó el manejo de hilos con el objetivo de permitir el procesamiento concurrente de múltiples archivos PDF dentro del programa.
Inicialmente, el sistema únicamente utilizaba un hilo secundario para detectar interrupciones por parte del usuario mediante el teclado. Posteriormente, se incorporó un modelo de procesamiento concurrente utilizando un pool de hilos administrado mediante \texttt{ExecutorService}.
Esto permitió que múltiples archivos PDF fueran procesados simultáneamente, aproximando el comportamiento del sistema a conceptos reales de concurrencia y planificación de tareas en Sistemas Operativos.
\vspace{-0.5cm}
\begin{figure}[H]
    \centering
    \caption{Pool de hilos para procesamiento concurrente}
    \label{code:pool_hilos}
    \begin{lstlisting}[language=Java]
ExecutorService pool =
    Executors.newFixedThreadPool(4);
List<Future<Alumno>> futures =
    new ArrayList<>();
for (Path pdf : stream) {
    System.out.println(
        "[MAIN] Enviando tarea: "
        + pdf.getFileName());
    ProcesadorPDFTask tarea =
        new ProcesadorPDFTask(pdf);
    futures.add(pool.submit(tarea));
}
    \end{lstlisting}
\end{figure}
Como se observa en la figura \ref{code:pool_hilos}, el sistema crea un pool fijo de cuatro hilos para administrar las tareas concurrentes. Cada archivo PDF es tratado como una tarea independiente y enviado al pool de ejecución mediante el método \texttt{submit()}.
Cada tarea es representada mediante la clase \texttt{ProcesadorPDFTask}, la cual implementa la interfaz \texttt{Callable<Alumno>} para permitir el procesamiento concurrente y el retorno de resultados.
\vspace{-0.5cm}
\begin{figure}[H]
    \centering
    \caption{Clase ProcesadorPDFTask}
    \label{code:procesador_task}
    \begin{lstlisting}[language=Java]
public class ProcesadorPDFTask
        implements Callable<Alumno> {
    private final Path pdf;
    public ProcesadorPDFTask(Path pdf) {
        this.pdf = pdf;
    }
    @Override
    public Alumno call() throws Exception {
        System.out.println(
            "[" + Thread.currentThread().getName()
            + "] Procesando: "
            + pdf.getFileName());
        try (PDDocument doc =
                 PDDocument.load(pdf.toFile())) {
            PDFHorarioStripper stripper =
                new PDFHorarioStripper();
            stripper.getText(doc);
            return new Alumno(
                stripper.getNombreAlumno(),
                stripper.getHorario()
            );
        }
    }
}
    \end{lstlisting}
\end{figure}
La clase \texttt{ProcesadorPDFTask} permite encapsular el procesamiento de cada archivo PDF dentro de una tarea concurrente independiente. Cada hilo del pool ejecuta distintas tareas simultáneamente, reutilizando recursos de ejecución de manera eficiente.
Asimismo, el sistema implementa tolerancia a fallos durante la ejecución concurrente. Si un archivo PDF presenta errores de lectura o corrupción, únicamente la tarea asociada falla, mientras las demás continúan ejecutándose normalmente.
\vspace{-0.5cm}
\begin{figure}[H]
    \centering
    \caption{Ejecución concurrente del sistema}
    \label{code:salida_hilos}
    \begin{verbatim}
[MAIN] Enviando tarea: HorarioAlfredo.pdf
[MAIN] Enviando tarea: HorarioAri.pdf
[pool-1-thread-1] Procesando: HorarioAlfredo.pdf
[pool-1-thread-2] Procesando: HorarioAri.pdf
    \end{verbatim}
\end{figure}
La implementación concurrente permitió mejorar el aprovechamiento de recursos del sistema y reducir el tiempo de procesamiento al ejecutar múltiples tareas de manera simultánea.

Para garantizar un cierre ordenado del pool de hilos, se implementó un mecanismo de terminación controlada. Una vez que todas las tareas fueron enviadas y sus resultados recolectados mediante \texttt{Future.get()}, se invoca \texttt{shutdown()} para indicar que no se aceptarán nuevas tareas. Posteriormente, \texttt{awaitTermination()} espera hasta que los hilos activos finalicen su ejecución. En caso de que el tiempo de espera se agote, se invoca \texttt{shutdownNow()} para forzar la interrupción de los hilos restantes.
\vspace{-0.5cm}
\begin{figure}[H]
    \centering
    \caption{Cierre controlado del pool de hilos}
    \label{code:shutdown_pool}
    \begin{lstlisting}[language=Java]
pool.shutdown();
try {
    if (!pool.awaitTermination(10, TimeUnit.SECONDS)) {
        pool.shutdownNow();
    }
} catch (InterruptedException e) {
    pool.shutdownNow();
    Thread.currentThread().interrupt();
}
    \end{lstlisting}
\end{figure}

Adicionalmente, el sistema incorpora un hilo secundario de tipo \textit{daemon} que permite al usuario cancelar la ejecución del programa durante la cuenta regresiva inicial presionando cualquier tecla. Al ser un hilo daemon, este termina automáticamente cuando el hilo principal concluye, sin necesidad de gestión explícita.

Para la detección de entrada por teclado se utilizó \texttt{Channels.newChannel(System.in)} en lugar de \texttt{System.in.read()} directamente. Esto se debe a que \texttt{System.in.read()} bloquea el hilo de forma no interrumpible, impidiendo que herramientas como Maven puedan terminarlo correctamente al finalizar el programa. El canal NIO, en cambio, lanza una \texttt{ClosedByInterruptException} al recibir una interrupción, permitiendo un cierre limpio.
\vspace{-0.5cm}
\begin{figure}[H]
    \centering
    \caption{Hilo daemon para detección de interrupción por teclado}
    \label{code:hilo_daemon}
    \begin{lstlisting}[language=Java]
Thread hiloInterrupcion = new Thread(() -> {
    try {
        ReadableByteChannel canal =
            Channels.newChannel(System.in);
        ByteBuffer buffer = ByteBuffer.allocate(1);
        if (canal.read(buffer) != -1) {
            System.out.println(
                "\n[!] Ejecucion cancelada.");
            System.exit(0);
        }
    } catch (ClosedByInterruptException e) {
        // Interrupcion limpia por Maven
    } catch (Exception e) {
        // Ignorado intencionalmente
    }
});
hiloInterrupcion.setDaemon(true);
hiloInterrupcion.start();
    \end{lstlisting}
\end{figure}

\subsection*{Sincronización y exclusión mutua}

Durante el diseño del sistema se evaluó la necesidad de mecanismos de sincronización para el acceso al archivo de salida. Dado que los hilos del pool únicamente realizan operaciones de \textbf{lectura} sobre sus respectivos archivos PDF y retornan un objeto \texttt{Alumno} como resultado, no existe acceso concurrente al archivo de salida en ningún momento de la ejecución.

La escritura sobre el archivo de salida ocurre exclusivamente en el hilo principal, una vez que \texttt{awaitTermination()} garantiza que todos los hilos del pool han finalizado. Por lo tanto, no se presenta ninguna condición de carrera y no se requiere el uso de \texttt{ReentrantLock} ni ningún otro mecanismo de exclusión mutua para proteger el recurso compartido.

Este análisis refleja una decisión de diseño consciente: incorporar sincronización innecesaria introduciría overhead sin beneficio real, por lo que se optó por un flujo secuencial para la fase de escritura.
